\chapter{Taxonomía de controles}
\label{ch:controles}

El modelo de madurez del capítulo anterior establece qué se evalúa y en qué escala, pero no especifica las medidas concretas que un sistema debe implementar para alcanzar cada nivel. Ese es el objeto de la taxonomía de controles: traducir cada nivel de cada subdimensión en un conjunto verificable de acciones. El capítulo define primero los atributos que caracterizan a un control, presenta luego el catálogo completo organizado por subdimensión, y cierra con la distribución de controles por nivel que resulta de la valoración de riesgo establecida en la sección~\ref{sec:valoracion-riesgo}.


\section{Estructura de los controles}

Cada control se caracteriza por cinco atributos:

\begin{enumerate}[leftmargin=2em]
  \item \textbf{Nivel de madurez} (N1 a N5): determina en qué nivel del modelo se requiere su implementación.
  \item \textbf{Tipo funcional}: preventivo, detectivo, correctivo o gobernanza, según el momento en que actúa respecto al evento de seguridad (ISO/IEC 27002:2022).
  \item \textbf{Fuente de evidencia}: configuración (archivos YAML, prompt del sistema), experimental (resultados de auditoría ofensiva) o declarativo (verificación por lista de chequeo).
  \item \textbf{Descripción}: acción específica que el control requiere.
  \item \textbf{Fuente normativa}: referencia al estándar o legislación que lo origina.
\end{enumerate}

Un control puede satisfacer más de un tipo funcional simultáneamente, consistente con el modelo de atributos de ISO/IEC 27002:2022. Por ejemplo, un escáner de salida que identifica y suprime contenido sensible actúa como control detectivo y correctivo en una misma operación.

Todos los controles de un mismo nivel tienen peso equivalente dentro de su subdimensión.

Un control transversal, aplicable a más de una subdimensión, puede citar fuentes normativas distintas en cada una de ellas. Esto responde a que su función protectora se articula contra vectores de amenaza diferentes según la etapa del flujo RAG en que actúa: el proxy de seguridad (C16) intercepta inyección de instrucciones en la entrada (LLM01), limita la exposición del corpus en la recuperación (LLM02) y previene la fuga del prompt del sistema en la salida (LLM07). El identificador y la descripción del control permanecen invariantes; solo el conjunto de fuentes normativas se acota a los vectores pertinentes a cada subdimensión.

El conjunto completo comprende 62 controles distribuidos en cinco subdimensiones y cinco niveles de madurez. Las secciones siguientes presentan los controles organizados por subdimensión.

% ══════════════════════════════════════════════
% D1: Protección de entrada (17 controles)
% ══════════════════════════════════════════════
\section{D1: Protección de entrada}
\label{sec:controles-d1}

La subdimensión D1 agrupa 17 controles distribuidos en cinco niveles de madurez: 8 en N1, 3 en N2, 2 en N3, 2 en N4 y 2 en N5. El predominio de controles preventivos refleja la naturaleza de esta subdimensión: interceptar la amenaza antes de que alcance la capa de inferencia.

\control{C01}{1}{Preventivo}{Configuración}
  {Limitar la longitud máxima del prompt de entrada, medida en la misma unidad que el modelo emplea para su ventana de contexto.}
  {LLM01:2025}

\control{C02}{1}{Correctivo}{Experimental}
  {Eliminar de la entrada los tokens especiales del modelo, los bytes nulos y los caracteres de control Unicode, antes de que el texto alcance cualquier componente de análisis.}
  {LLM01:2025}

\control{C06}{1}{Preventivo}{Configuración}
  {Segregar e identificar el contenido recuperado del corpus mediante delimitadores explícitos que lo separen de la consulta del usuario en el prompt aumentado.}
  {LLM01:2025, LLM02:2025}

\control{C13}{1}{Preventivo}{Configuración}
  {Limitar el volumen total de contexto inyectado al LLM restringiendo la cantidad de fragmentos recuperados por consulta.}
  {LLM01:2025, LLM02:2025}

\control{C15}{1}{Preventivo}{Configuración}
  {Restringir el rol, las capacidades y las limitaciones del modelo mediante instrucciones de sistema declaradas externamente al código.}
  {LLM01:2025, LLM02:2025}

\control{C16}{1}{Preventivo}{Configuración}
  {Interponer un proxy de seguridad desplegado como componente independiente del generador, que intercepte la entrada y la salida del flujo de inferencia.}
  {LLM01:2025, LLM02:2025}

\control{C21}{1}{Preventivo}{Configuración}
  {Externalizar la configuración de tiempo de ejecución, incluidas las credenciales y los umbrales, como parámetros declarados fuera del código fuente.}
  {LLM01:2025, LLM02:2025}

\control{C30}{1}{Gobernanza}{Declarativo}
  {Verificar la procedencia e integridad del modelo mediante \textit{checksums} antes del despliegue.}
  {Ley 21.663, LLM08:2025}

\control{C03}{2}{Preventivo}{Experimental}
  {Bloquear las entradas clasificadas como intento de inyección de instrucciones por un escáner que asigne una puntuación de riesgo contrastada contra un umbral configurado.}
  {LLM01:2025}

\control{C04}{2}{Preventivo}{Experimental}
  {Bloquear las consultas cuya distancia semántica mínima al corpus supere el umbral configurado, rechazando aquellas ajenas al dominio autorizado.}
  {LLM01:2025}

\control{C05}{2}{Detectivo}{Experimental}
  {Extraer de la respuesta del modelo únicamente el texto generado, sin incorporar el prompt del sistema ni sus restricciones al contenido retornado al usuario.}
  {LLM01:2025, LLM07:2025}

\control{C31}{3}{Detectivo}{Experimental}
  {Registrar un motivo de bloqueo estructurado para cada solicitud rechazada, identificando el componente que originó el rechazo.}
  {LLM01:2025}

\control{C32}{3}{Preventivo}{Configuración}
  {Verificar mediante función de resumen criptográfico que el prompt del sistema y sus restricciones no han sido modificados desde el arranque.}
  {LLM01:2025, LLM07:2025}

\control{C33}{4}{Preventivo}{Declarativo}
  {Ajustar dinámicamente los umbrales de detección de inyección según las tasas de falsos positivos observadas.}
  {LLM01:2025}

\control{C34}{4}{Detectivo}{Declarativo}
  {Medir y reportar periódicamente la tasa de prompts rechazados por categoría de ataque.}
  {LLM01:2025}

\control{C35}{5}{Preventivo}{Declarativo}
  {Actualizar los patrones de detección ante nuevas técnicas de inyección documentadas.}
  {LLM01:2025}

\control{C36}{5}{Preventivo}{Declarativo}
  {Revisar y ajustar el prompt del sistema ante cambios en el modelo base.}
  {LLM01:2025}

% ══════════════════════════════════════════════
% D2: Control de contexto RAG (20 controles)
% ══════════════════════════════════════════════
\section{D2: Control de contexto RAG}
\label{sec:controles-d2}

La subdimensión D2 agrupa 20 controles orientados a regular el acceso, la exposición y la integridad de la base de datos vectorial y el corpus de conocimiento. La distribución por nivel refleja una progresión desde controles básicos de particionamiento y sanitización (N1) hasta mecanismos de auditoría avanzada y mejora continua (N5).

\control{C06}{1}{Preventivo}{Configuración}
  {Segregar e identificar el contenido recuperado del corpus mediante delimitadores explícitos que lo separen de la consulta del usuario en el prompt aumentado.}
  {LLM01:2025, LLM02:2025}

\control{C07}{1}{Preventivo}{Configuración}
  {Restringir el modelo al uso exclusivo del contexto recuperado, prohibiendo mediante instrucción de sistema la generación de información no fundamentada en él.}
  {LLM02:2025}

\control{C10}{1}{Preventivo}{Configuración}
  {Controlar la cantidad de fragmentos recuperados por consulta para limitar la exposición del corpus.}
  {LLM02:2025}

\control{C13}{1}{Preventivo}{Configuración}
  {Limitar el volumen total de contexto inyectado al LLM restringiendo la cantidad de fragmentos recuperados por consulta.}
  {LLM02:2025}

\control{C15}{1}{Preventivo}{Configuración}
  {Restringir el rol, las capacidades y las limitaciones del modelo mediante instrucciones de sistema declaradas externamente al código.}
  {LLM01:2025, LLM02:2025}

\control{C16}{1}{Preventivo}{Configuración}
  {Interponer un proxy de seguridad desplegado como componente independiente del generador, que intercepte la entrada y la salida del flujo de inferencia.}
  {LLM01:2025, LLM02:2025}

\control{C21}{1}{Preventivo}{Configuración}
  {Externalizar la configuración de tiempo de ejecución, incluidas las credenciales y los umbrales, como parámetros declarados fuera del código fuente.}
  {LLM01:2025, LLM02:2025}

\control{C30}{1}{Gobernanza}{Declarativo}
  {Verificar la procedencia e integridad del modelo mediante \textit{checksums} antes del despliegue.}
  {Ley 21.663, LLM08:2025}

\control{C08}{2}{Preventivo}{Configuración}
  {Clasificar cada fragmento del corpus según su nivel de sensibilidad al momento de la indexación, y restringir la recuperación a los niveles autorizados.}
  {LLM02:2025, LLM08:2025}

\control{C09}{2}{Preventivo/Detectivo}{Experimental}
  {Verificar la integridad de los fragmentos recuperados contra una línea base establecida, antes de inyectarlos al contexto de inferencia.}
  {LLM02:2025, LLM08:2025}

\control{C11}{2}{Preventivo}{Experimental}
  {Bloquear las consultas cuya distancia semántica mínima a los fragmentos del corpus supere el umbral de coherencia configurado, rechazando aquellas incoherentes o codificadas.}
  {LLM02:2025, LLM08:2025}

\control{C14}{2}{Preventivo}{Experimental}
  {Excluir los metadatos de los documentos en la etapa de recuperación, de modo que solo el texto de los fragmentos alcance el contexto de inferencia.}
  {LLM02:2025}

\control{C17}{2}{Preventivo/Detectivo}{Experimental}
  {Bloquear las respuestas cuyo solapamiento con el contexto recuperado caiga por debajo del umbral de fundamentación configurado.}
  {LLM02:2025, LLM07:2025}

\control{C22}{2}{Preventivo/Detectivo}{Declarativo}
  {Controlar y auditar el acceso por sesión a los documentos del \textit{vector store}.}
  {LLM02:2025}

\control{C12}{3}{Detectivo}{Experimental}
  {Detectar y bloquear las consultas indicativas de reconocimiento o enumeración sistemática del corpus.}
  {LLM02:2025, LLM08:2025}

\control{C62}{3}{Detectivo}{Configuración}
  {Mantener un registro versionado y persistente de cada evento de indexación del corpus, que permita identificar qué documentos estaban activos en cada sesión y rastrear los cambios entre versiones.}
  {LLM02:2025, LLM08:2025}

\control{C37}{4}{Detectivo}{Declarativo}
  {Medir y reportar periódicamente la tasa de recuperación de fragmentos sensibles ante consultas de prueba.}
  {LLM02:2025}

\control{C38}{4}{Preventivo}{Declarativo}
  {Ajustar los parámetros de \textit{retrieval} según las métricas de exposición observadas.}
  {LLM02:2025}

\control{C39}{5}{Preventivo}{Declarativo}
  {Actualizar el particionamiento del \textit{vector store} ante cambios en la clasificación de sensibilidad del corpus.}
  {LLM02:2025, LLM08:2025}

\control{C40}{5}{Preventivo}{Declarativo}
  {Revisar proactivamente los fragmentos indexados ante nuevas amenazas de exfiltración documentadas.}
  {LLM02:2025, LLM08:2025}

% ══════════════════════════════════════════════
% D3: Gestión de salida (10 controles)
% ══════════════════════════════════════════════
\section{D3: Gestión de salida}
\label{sec:controles-d3}

La subdimensión D3 comprende 9 controles que filtran y regulan las respuestas generadas por el modelo antes de entregarlas al usuario, constituyendo la última línea de defensa contra la exfiltración de información sensible y la fuga del prompt del sistema.

\control{C16}{1}{Preventivo}{Configuración}
  {Interponer un proxy de seguridad desplegado como componente independiente del generador, que intercepte la entrada y la salida del flujo de inferencia.}
  {LLM01:2025, LLM02:2025, LLM07:2025}

\control{C17}{2}{Preventivo/Detectivo}{Experimental}
  {Bloquear las respuestas cuyo solapamiento con el contexto recuperado caiga por debajo del umbral de fundamentación configurado.}
  {LLM02:2025, LLM07:2025}

\control{C18}{2}{Correctivo}{Experimental}
  {Eliminar de la salida del modelo los caracteres de control y los tokens especiales, antes de someterla a los escáneres de contenido y de retornarla al usuario.}
  {LLM01:2025, LLM07:2025}

\control{C19}{2}{Detectivo/Correctivo}{Experimental}
  {Bloquear las respuestas que coincidan con patrones deterministas de contenido sensible: marcadores de confidencialidad e identificadores personales estructurados.}
  {LLM02:2025, LLM07:2025}

\control{C20}{3}{Detectivo}{Experimental}
  {Detectar y puntuar mediante análisis semántico las respuestas que contengan información sensible no capturada por los patrones deterministas.}
  {LLM02:2025, LLM07:2025}

\control{C41}{4}{Detectivo}{Declarativo}
  {Ajustar dinámicamente los umbrales de sanitización según las métricas de falsos positivos observadas en la salida.}
  {LLM02:2025, LLM07:2025}

\control{C42}{4}{Detectivo}{Declarativo}
  {Medir y reportar la tasa de supresión de fragmentos confidenciales por tipo de ataque.}
  {LLM02:2025, LLM07:2025}

\control{C43}{5}{Detectivo}{Declarativo}
  {Revisar y actualizar los patrones de detección de contenido sensible en la salida ante nuevas técnicas de exfiltración.}
  {LLM02:2025, LLM07:2025}

\control{C44}{5}{Detectivo}{Declarativo}
  {Incorporar retroalimentación de incidentes reales al conjunto de patrones de sanitización.}
  {LLM02:2025, LLM07:2025}

% ══════════════════════════════════════════════
% D4: Monitoreo y trazabilidad (9 controles)
% ══════════════════════════════════════════════
\section{D4: Monitoreo y trazabilidad}
\label{sec:controles-d4}

La subdimensión D4 agrupa 9 controles de naturaleza observacional, distribuidos a partir de N2. D4 no define controles de nivel N1 porque el monitoreo mínimo viable requiere que el sistema disponga de algún mecanismo activo de registro, condición que no se cumple en un sistema sin capa defensiva.

\control{C26}{2}{Gobernanza}{Declarativo}
  {Definir y documentar la política de umbrales de detección y bloqueo con anterioridad a la operación.}
  {Ley 21.663}

\control{C27}{2}{Gobernanza}{Declarativo}
  {Establecer un procedimiento formal de notificación de incidentes de seguridad ante las autoridades competentes.}
  {Ley 21.663}

\control{C28}{2}{Gobernanza}{Declarativo}
  {Definir la política de retención y eliminación de datos de sesión conforme a la normativa aplicable.}
  {Ley 21.719}

\control{C23}{3}{Detectivo}{Experimental}
  {Registrar cada decisión de bloqueo en la entrada y en la salida, identificando el componente que la originó.}
  {LLM01:2025, LLM02:2025, LLM07:2025, LLM08:2025}

\control{C25}{3}{Detectivo}{Experimental}
  {Registrar la latencia y los metadatos de cada solicitud en un formato que permita el análisis posterior del comportamiento del flujo de inferencia.}
  {LLM01:2025, LLM02:2025, LLM07:2025, LLM08:2025}

\control{C29}{3}{Gobernanza}{Declarativo}
  {Documentar y ejecutar un plan de pruebas de seguridad ofensiva con anterioridad a cada despliegue.}
  {Ley 21.663}

\control{C24}{4}{Detectivo}{Experimental}
  {Monitorear los patrones de consulta anómalos que sugieran reconocimiento sistemático del corpus a lo largo de múltiples sesiones.}
  {LLM02:2025, LLM08:2025}

\control{C45}{5}{Detectivo}{Declarativo}
  {Revisar y actualizar los patrones de detección de anomalías ante nuevas amenazas documentadas.}
  {LLM01:2025, LLM02:2025, LLM07:2025, LLM08:2025}

\control{C46}{5}{Correctivo}{Declarativo}
  {Incorporar retroalimentación de incidentes al sistema de alertas de forma sistemática y periódica.}
  {LLM01:2025, LLM02:2025, LLM07:2025, LLM08:2025}



% ══════════════════════════════════════════════
% D5: Gobernanza organizacional (15 controles)
% ══════════════════════════════════════════════
\section{D5: Gobernanza organizacional}
\label{sec:controles-d5}

La subdimensión D5 agrupa 15 controles de tipo exclusivamente Gobernanza con fuente de evidencia Declarativa. Estos controles capturan las obligaciones institucionales derivadas de la Ley 21.663, la Ley 21.719, el NIST AI RMF y el Proyecto de Ley de IA.

\control{C47}{1}{Gobernanza}{Declarativo}
  {Designar formalmente un responsable de la seguridad del sistema RAG con atribuciones documentadas.}
  {Ley 21.663}

\control{C48}{1}{Gobernanza}{Declarativo}
  {Declarar el sistema RAG como activo de información sujeto a las obligaciones de la normativa vigente.}
  {Ley 21.663}

\control{C49}{1}{Gobernanza}{Declarativo}
  {Establecer una política de uso aceptable del sistema RAG accesible a todos los usuarios.}
  {Ley 21.719}

\control{C61}{1}{Gobernanza}{Declarativo}
  {Mantener un inventario actualizado de los activos del sistema RAG, incluyendo modelo, corpus, base vectorial y componentes de seguridad, conforme a lo definido en la sección~\ref{sec:activos}.}
  {Ley 21.663}

\control{C50}{2}{Gobernanza}{Declarativo}
  {Adoptar y adaptar la matriz RACI del marco LLM-Sec, asignando los roles genéricos definidos a personas o cargos concretos de la organización.}
  {Ley 21.663}

\control{C51}{2}{Gobernanza}{Declarativo}
  {Establecer un procedimiento formal de gestión de accesos que regule altas, bajas y modificaciones de permisos sobre el sistema.}
  {Ley 21.719}

\control{C52}{2}{Gobernanza}{Declarativo}
  {Documentar el inventario de datos personales procesados por el sistema RAG conforme a las obligaciones del responsable de tratamiento.}
  {Ley 21.719}

\control{C53}{3}{Gobernanza}{Declarativo}
  {Establecer un programa formal de capacitación en seguridad de sistemas LLM dirigido a los roles definidos en la matriz RACI.}
  {NIST AI RMF}

\control{C54}{3}{Gobernanza}{Declarativo}
  {Definir y ejecutar un ciclo de revisión periódica del marco LLM-Sec con frecuencia mínima anual.}
  {NIST AI RMF}

\control{C55}{3}{Gobernanza}{Declarativo}
  {Documentar y comunicar formalmente los riesgos residuales del sistema RAG a la dirección de la organización.}
  {NIST AI RMF}

\control{C60}{3}{Gobernanza}{Declarativo}
  {Realizar y documentar una evaluación formal de riesgos del sistema RAG con frecuencia mínima anual o ante cambios significativos en el modelo o el corpus.}
  {NIST AI RMF}

\control{C56}{4}{Gobernanza}{Declarativo}
  {Establecer métricas de desempeño de la gobernanza de seguridad y reportarlas periódicamente a la dirección.}
  {Proyecto Ley IA}

\control{C57}{4}{Gobernanza}{Declarativo}
  {Implementar un proceso formal de gestión de riesgos de IA que clasifique el sistema RAG según su nivel de criticidad.}
  {Proyecto Ley IA}

\control{C58}{5}{Gobernanza}{Declarativo}
  {Revisar y actualizar proactivamente el marco LLM-Sec ante cambios normativos, nuevas amenazas o modificaciones del modelo base.}
  {Proyecto Ley IA}

\control{C59}{5}{Gobernanza}{Declarativo}
  {Incorporar lecciones aprendidas de incidentes de seguridad al ciclo de mejora continua del marco de gobernanza.}
  {Proyecto Ley IA}

% ══════════════════════════════════════════════
% Resumen de distribución
% ══════════════════════════════════════════════
\section{Distribución de controles}
\label{sec:distribucion}

El marco define 62 controles con identificador único (C01 a C62). Dado que un control puede proteger activos correspondientes a más de una subdimensión, un mismo identificador aparece en cada una de las subdimensiones donde su función protectora es aplicable. El control C16 (proxy de seguridad), por ejemplo, opera simultáneamente sobre la entrada, el contexto recuperado y la salida, por lo que se lista en D1, D2 y D3.

La Tabla~\ref{tab:distribucion} contabiliza aplicaciones de control por subdimensión y nivel de madurez, no controles únicos. La suma total es por tanto superior a 62: los 62 controles únicos generan 70 aplicaciones, diferencia que corresponde a los siete controles transversales del marco (C06, C13, C15, C16, C17, C21 y C30).

\begin{fichabox}{Distribución de controles: subdimensión $\times$ nivel}
  \label{tab:distribucion}
  \begin{tabularx}{\textwidth}{@{}l C{1.5cm} C{1.5cm} C{1.5cm} C{1.5cm} C{1.5cm} C{2.2cm}@{}}
    \toprule
    \textbf{Subdim.} & \textbf{N1} & \textbf{N2} & \textbf{N3} & \textbf{N4} & \textbf{N5} & \textbf{Aplicaciones} \\
    \midrule
    D1 & 8 & 3 & 2 & 2 & 2 & 17 \\
    D2 & 8 & 6 & 2 & 2 & 2 & 20 \\
    D3 & 1 & 3 & 1 & 2 & 2 &  9 \\
    D4 & 0 & 3 & 3 & 1 & 2 &  9 \\
    D5 & 4 & 3 & 4 & 2 & 2 & 15 \\
    \midrule
    \textbf{Total} & \textbf{21} & \textbf{18} & \textbf{12} & \textbf{9} & \textbf{10} & \textbf{70} \\
    \bottomrule
  \end{tabularx}
\end{fichabox}